\setcounter{section}{26}

  \zwischenueberschrift{Soal latihan}

\inputaufgabe
{}
{

Misalkan bidang real dilengkapi dengan
\definitionsverweis {struktur Euklides}{}{}
dan
\definitionsverweis {koneksi Levi--Civita}{}{}
yang bersesuaian. Hitung

\mavergleichskette
{\vergleichskette
{ \nabla_FG
}
{ = }{
}
{  }{
}
{    }{
}
{   }{
}
}
{}{}{}
untuk medan-medan vektor

\mavergleichskettedisp
{\vergleichskette
{ F(x,y)
}
{ =} { \left(  x^2-y^3  , \,   xy+y^2                   \right)
}
{ } {
}
{ } {
}
{ } {
}
}
{}{}{}
dan

\mavergleichskettedisp
{\vergleichskette
{ G(x,y)
}
{ =} { \left( e^{xy^2}   , \,   x^3-x^2y^5                   \right)
}
{ } {
}
{ } {
}
{ } {
}
}
{}{}{.}

}
{} {}

\inputaufgabe
{}
{

Misalkan

\mavergleichskette
{\vergleichskette
{ I
}
{ \subseteq }{ \R
}
{  }{
}
{    }{
}
{   }{
}
}
{}{}{}
suatu
\definitionsverweis {interval real}{}{}
dan misalkan

\mavergleichskettedisp
{\vergleichskette
{ g(t)
}
{ =} { 3+t^2
}
{ } {
}
{ } {
}
{ } {
}
}
{}{}{,}
yang kita tafsirkan sebagai suatu
\definitionsverweis {metrik Riemann}{}{}
pada $I$. Tentukan
\definitionsverweis {turunan vertikal}{}{}
\zusatzklammer {terhadap
\definitionsverweis {koneksi Levi--Civita}{}{}
yang bersesuaian} {} {}

\mathl{\nabla_\partial  \left(  t^3-  \sin  \left( t^2 \right)   \right)}{.}

}
{} {}

\inputaufgabegibtloesung
{}
{

Misalkan

\mavergleichskette
{\vergleichskette
{ I
}
{ \subseteq }{ \R
}
{  }{
}
{    }{
}
{   }{
}
}
{}{}{}
suatu
\definitionsverweis {interval real}{}{}
dan misalkan
\maabbdisp {g} { I } {\R_+
} {}
suatu fungsi positif diferensiabel yang kita tafsirkan sebagai
\definitionsverweis {metrik Riemann}{}{}
pada $I$. Tentukan
\definitionsverweis {penampang horizontal}{}{}
terhadap
\definitionsverweis {koneksi Levi--Civita}{}{}
yang bersesuaian pada $I \times \R$
\zusatzklammer {lihat
[[Intervall/Riemannsche Struktur/Levi-Civita-Zusammenhang/Beispiel|Contoh 26.4]]} {} {.}

}
{} {}

\inputaufgabe
{}
{

Misalkan $M$ suatu
\definitionsverweis {manifold diferensiabel}{}{}
sedemikian sehingga
\definitionsverweis {bundel tangen}{}{}
$TM$ trivial, dan misalkan pada $M$, melalui suatu trivialisasi terdiferensialkan

\mavergleichskette
{\vergleichskette
{ TM
}
{ \cong }{ M \times \R^n
}
{  }{
}
{    }{
}
{   }{
}
}
{}{}{}
\zusatzklammer {dengan menggunakan
\definitionsverweis {hasil kali skalar standar}{}{}
pada $\R^n$} {} {}
diberikan suatu
\definitionsverweis {struktur Riemann}{}{}
\zusatzklammer {lihat
[[Mannigfaltigkeit/Tangentialbündel trivial/Riemannsche Struktur/Aufgabe|Soal 16.2]]} {} {.}
Tunjukkan bahwa
\definitionsverweis {koneksi Levi--Civita}{}{}
pada $M$ adalah
\definitionsverweis {koneksi trivial}{}{}
tersebut.

}
{} {}

\inputaufgabe
{}
{

Misalkan ${\mathbb H}$ adalah
\definitionsverweis {setengah bidang Poincaré}{}{}
dengan
\definitionsverweis {struktur Riemann}{}{}
dan
\definitionsverweis {koneksi Levi--Civita}{}{}
yang bersesuaian. Tentukan
\definitionsverweis {turunan vertikal}{}{}
\mathl{\nabla_{\partial_1 + \partial_2 } (G)}{} dari medan vektor

\mavergleichskettedisp
{\vergleichskette
{ G(x,y)
}
{ =} { \left(  x^5-y^3+2xy  , \,   4 x^2+ y^3                   \right)
}
{ } {
}
{ } {
}
{ } {
}
}
{}{}{.}

}
{} {}

\inputaufgabegibtloesung
{}
{

Misalkan diberikan suatu
\definitionsverweis {koneksi linear}{}{}
$\nabla$ pada
\definitionsverweis {bundel tangen}{}{}
$TX$ dari suatu
\definitionsverweis {manifold diferensiabel}{}{}
$X$. Andaikan terdapat
\definitionsverweis {penampang-penampang basis}{}{}

\mathl{V_1 , \ldots , V_n}{} dengan

\mavergleichskettedisp
{\vergleichskette
{ \nabla_{V_i } V_j -  \nabla_{V_j } V_i
}
{ =} { [V_i,V_j]
}
{ } {
}
{ } {
}
{ } {
}
}
{}{}{}
untuk semua $i,j$. Tunjukkan bahwa $\nabla$
\definitionsverweis {bebas torsi}{}{.}

}
{} {}

\inputaufgabe
{}
{

Tunjukkan bahwa suatu
\definitionsverweis {koneksi linear}{}{}
pada
\definitionsverweis {bundel tangen}{}{}
dari interval terbuka

\mavergleichskette
{\vergleichskette
{ I
}
{ \subseteq }{ \R
}
{  }{
}
{    }{
}
{   }{
}
}
{}{}{}
bersifat
\definitionsverweis {bebas torsi}{}{.}

}
{} {}

\inputaufgabe
{}
{

Misalkan

\mavergleichskette
{\vergleichskette
{ U
}
{ \subseteq }{ \R^n
}
{  }{
}
{    }{
}
{   }{
}
}
{}{}{}
\definitionsverweis {terbuka}{}{}
dan padanya diberikan suatu
\definitionsverweis {struktur Riemann}{}{}
melalui fungsi-fungsi $\left(  g_{ij}  \right)$, dengan
\definitionsverweis {matriks invers}{}{}

\mathl{\left(  g^{k l}  \right)}{.} Misalkan $\nabla$ adalah koneksi yang diberikan oleh
\definitionsverweis {simbol-simbol Christoffel}{}{}

\mavergleichskettedisp
{\vergleichskette
{ \Gamma_{ij}^k
}
{ =} { { \frac{   1  }{  2 } } \sum_{l = 1}^n g^{kl}(\partial_ig_{jl}+\partial_jg_{il}-\partial_lg_{ij})
}
{ } {
}
{ } {
}
{ } {
}
}
{}{}{}
sebagai
\definitionsverweis {koneksi linear}{}{,}
yang dicirikan oleh

\mavergleichskettedisp
{\vergleichskette
{ \nabla_{\partial_i } \partial_j
}
{ =} { \sum_{k = 1}^n\Gamma_{ij}^k\,\partial_k
}
{ } {
}
{ } {
}
{ } {
}
}
{}{}{.}
Tunjukkan dalam dua langkah bahwa $\nabla$ memenuhi rumus Koszul dari
[[Riemannsche Mannigfaltigkeit/Tangentialbündel/Linearer Zusammenhang/Metrisch und torsionsfrei/Koszul-Eigenschaft/Eindeutigkeit/Fakt|Lemma 26.10]].
\aufzaehlungzwei {Untuk medan-medan basis $\partial_i$.
} {Untuk medan-medan vektor sembarang $V,W,Z$.
}

}
{} {}

\inputaufgabegibtloesung
{}
{

Misalkan $M$ suatu
\definitionsverweis {manifold Riemann}{}{}
dan $\nabla$ suatu
\definitionsverweis {koneksi linear}{}{}
pada $TM$. Andaikan bahwa $\nabla$, secara lokal untuk medan-medan basis $\partial_1  , \ldots , \partial_n$, memenuhi sifat

\mavergleichskettedisp
{\vergleichskette
{ \partial_i \left\langle  \partial_j  ,  \partial_k   \right\rangle
}
{ =} { \left\langle \nabla_{\partial_i } \partial_j   ,  \partial_k   \right\rangle  + \left\langle  \partial_j  ,  \nabla_{\partial_i } \partial_k    \right\rangle
}
{ } {
}
{ } {
}
{ } {
}
}
{}{}{}
untuk semua $i,j,k$. Tunjukkan bahwa $\nabla$
\definitionsverweis {metrik}{}{.}

}
{} {}

Konsep koneksi metrik tidak hanya terdapat pada bundel tangen suatu manifold Riemann, tetapi secara lebih umum juga pada bundel vektor Riemann.

Misalkan
\maabb {p} { E } { M
} {}
suatu
\definitionsverweis {bundel vektor Riemann}{}{}
di atas suatu
\definitionsverweis {manifold diferensiabel}{}{}
$M$. Suatu
\definitionsverweis {koneksi linear}{}{}
pada $E$ disebut
\definitionswort {metrik}{,}
jika pada setiap himpunan terbuka berlaku

\mavergleichskettedisp
{\vergleichskette
{ D_V \left\langle  s  , t  \right\rangle
}
{ =} { \left\langle \nabla_V s , t  \right\rangle  + \left\langle  s  ,  \nabla_Vt  \right\rangle
}
{ } {
}
{ } {
}
{ } {
}
}
{}{}{}
untuk suatu
\definitionsverweis {medan vektor}{}{}
$V$ yang
\definitionsverweis {terdiferensialkan secara kontinu}{}{}
dan penampang-penampang $s,t$ dari $E$ yang terdiferensialkan secara kontinu.

\inputaufgabe
{}
{

Misalkan

\mavergleichskette
{\vergleichskette
{ I
}
{ \subseteq }{ \R
}
{  }{
}
{    }{
}
{   }{
}
}
{}{}{}
suatu
\definitionsverweis {interval terbuka}{}{}
dan misalkan
\maabbdisp {p} {E = I \times \R^n } { I
} {}
bundel vektor trivial di atas $I$, yang kita lengkapi dengan
\definitionsverweis {hasil kali skalar standar}{}{}
dan kita pandang sebagai
\definitionsverweis {bundel vektor Riemann}{}{.}
Tunjukkan bahwa suatu
\definitionsverweis {koneksi linear}{}{}
$\nabla$ pada $E$ bersifat
\definitionsverweis {metrik}{}{}
jika dan hanya jika
\definitionsverweis {simbol-simbol Christoffel}{}{}
memenuhi syarat
\zusatzklammer {semuanya mengacu pada satu-satunya arah turunan $\partial$} {} {}

\mavergleichskettedisp
{\vergleichskette
{ \Gamma_j^k
}
{ =} { - \Gamma_k^j
}
{ } {
}
{ } {
}
{ } {
}
}
{}{}{}
untuk

\mavergleichskette
{\vergleichskette
{ 1
}
{ \leq }{ j,k
}
{ \leq }{ n
}
{    }{
}
{   }{
}
}
{}{}{.}

}
{} {}

\inputaufgabe
{}
{

Misalkan
\maabb {p} { E } { M
} {}
suatu
\definitionsverweis {bundel vektor Riemann}{}{}
pada suatu
\definitionsverweis {manifold diferensiabel}{}{}
$M$, dan misalkan $\nabla$ suatu
\definitionsverweis {koneksi linear}{}{}
pada $E$. Andaikan bahwa $\nabla$, secara lokal untuk
\definitionsverweis {medan-medan vektor standar}{}{}
$\partial_1  , \ldots , \partial_n$ dan
\definitionsverweis {penampang-penampang basis}{}{}

\mathl{s_1  , \ldots , s_r}{} dari $E$, memenuhi sifat

\mavergleichskettedisp
{\vergleichskette
{ \partial_i \left\langle  s_j  , s_k   \right\rangle
}
{ =} { \left\langle \nabla_{\partial_i } s_j   ,  s_k   \right\rangle  + \left\langle s_j  ,  \nabla_{\partial_i } s_k    \right\rangle
}
{ } {
}
{ } {
}
{ } {
}
}
{}{}{}
untuk semua $i,j,k$. Tunjukkan bahwa $\nabla$
\definitionsverweis {metrik}{}{.}

}
{} {}

Untuk soal berikut, bandingkan dengan
[[Eingebettete Mannigfaltigkeit/Inverse Karte/Tangentialbündel/Partielle Ableitung/Aufgabe|Soal 10.16]].

\inputaufgabe
{}
{

Misalkan

\mavergleichskette
{\vergleichskette
{ Y
}
{ \subseteq }{ \R^n
}
{  }{
}
{    }{
}
{   }{
}
}
{}{}{}
suatu
\definitionsverweis {submanifold tertutup}{}{}
dan
\maabb {\varphi} { V } { U \subseteq Y
} {}
suatu parametrisasi $C^2$-difeomorfik dari suatu himpunan bagian terbuka

\mavergleichskette
{\vergleichskette
{ U
}
{ \subseteq }{ Y
}
{  }{
}
{    }{
}
{   }{
}
}
{}{}{}
oleh suatu himpunan bagian terbuka

\mavergleichskette
{\vergleichskette
{ V
}
{ \subseteq }{ \R^d
}
{  }{
}
{    }{
}
{   }{
}
}
{}{}{,}
di mana $\varphi$ juga kita pandang sebagai pemetaan ke $\R^n$. Tunjukkan bahwa terdapat diagram komutatif

\mathdisp {\begin{matrix}  V  & \stackrel{ \partial_i }{\longrightarrow} &  TV & \stackrel{ T( \partial_j) }{\longrightarrow} & TTV & \\ \!\!\!\!\! \,\, \, \varphi \downarrow & & \!\!\!\!\! T(\varphi) \downarrow & &  \downarrow TT(\varphi) \!\!\!\!\! & & \\  U  & \stackrel{  }{\longrightarrow} &  TU & \stackrel{  }{\longrightarrow} & TTU & \!\!\!\!\!  \\ \end{matrix}} {  }

seperti di atas. Bagaimana pemetaan-pemetaan pada diagonal dapat dideskripsikan?

}
{} {}

  \zwischenueberschrift{Tugas untuk dikumpulkan}

\inputaufgabe
{2}
{

Misalkan

\mavergleichskette
{\vergleichskette
{ I
}
{ \subseteq }{ \R
}
{  }{
}
{    }{
}
{   }{
}
}
{}{}{}
suatu
\definitionsverweis {interval real}{}{}
dan misalkan

\mavergleichskettedisp
{\vergleichskette
{ g(t)
}
{ =} { 1+t^4
}
{ } {
}
{ } {
}
{ } {
}
}
{}{}{,}
yang kita tafsirkan sebagai suatu
\definitionsverweis {metrik Riemann}{}{}
pada $I$. Tentukan
\definitionsverweis {turunan vertikal}{}{}
\zusatzklammer {terhadap
\definitionsverweis {koneksi Levi--Civita}{}{}
yang bersesuaian} {} {}

\mathl{\nabla_\partial \left(  3e^{-t^2} + t  \cos \left(  t  \right) - t^5  \right)}{.}

}
{} {}

\inputaufgabe
{2}
{

Hitung $\Gamma^1_{12},  \Gamma^2_{12},  \Gamma^1_{22},  \Gamma^2_{22}$ dalam
[[Halbebene/Poincaré/Riemannsche Geometrie/Beispiel|Contoh 18.8]].

}
{} {}

\inputaufgabe
{3}
{

Misalkan ${\mathbb H}$ adalah
\definitionsverweis {setengah bidang Poincaré}{}{}
dengan
\definitionsverweis {struktur Riemann}{}{}
dan
\definitionsverweis {koneksi Levi--Civita}{}{}
yang bersesuaian. Tentukan
\definitionsverweis {turunan vertikal}{}{}
\mathl{\nabla_{xy\partial_1 -2y^3 \partial_2 } (G)}{} dari medan vektor

\mavergleichskettedisp
{\vergleichskette
{ G(x,y)
}
{ =} { \left(  x^4-2y^2  , \,   3 x^3-  \ln  \left(  1+x^2y^2 \right)                   \right)
}
{ } {
}
{ } {
}
{ } {
}
}
{}{}{.}

}
{} {}

\inputaufgabe
{2}
{

Misalkan $\nabla$ suatu
\definitionsverweis {koneksi bebas torsi}{}{}
\definitionsverweis {linear}{}{}
pada suatu himpunan terbuka

\mavergleichskette
{\vergleichskette
{ U
}
{ \subseteq }{ \R^n
}
{  }{
}
{    }{
}
{   }{
}
}
{}{}{.}
Oleh berapa banyak
\definitionsverweis {simbol Christoffel}{}{}
yang saling bebas koneksi itu ditentukan
\zusatzklammer {sebagai fungsi dari $n$} {} {}?

}
{} {}

 [[Kategorie:Latexseite]]
